% §7 Decision Framework
% Status: Pending data

\section{Decision Framework}
\label{sec:framework}

\subsection{Workload Characterization}

Before a recommendation can be made, five workload characteristics must be measured:
memory regularity ($\in [0,1]$), arithmetic intensity (FLOP/byte), kernel duration ($\mu$s),
control flow divergence ($\in [0,1]$), and data structure type.

\subsection{Decision Logic}

TODO: describe the rule-based logic from analysis/decision\_framework.py.
Insert Fig~\ref{fig:framework_flowchart}.

\begin{figure}[ht]
    \centering
    % \includegraphics[width=0.8\linewidth]{figures/fig20_decision_framework.pdf}
    \caption{Decision framework flowchart. Input: workload profile. Output: abstraction strategy.}
    \label{fig:framework_flowchart}
\end{figure}

\subsection{Empirically Derived Thresholds}

The following thresholds are determined from the experimental campaign:

\begin{table}[ht]
\centering
\caption{Empirically derived decision thresholds. TODO: fill in after experiments.}
\begin{tabular}{lrr}
\toprule
Threshold & Hypothetical & Empirical \\
\midrule
Memory regularity cutoff & 0.8 & TODO \\
Minimum kernel duration & 100 $\mu$s & TODO \\
PPC success threshold & 0.80 & TODO \\
Control divergence limit & 0.3 & TODO \\
\bottomrule
\end{tabular}
\label{tab:thresholds}
\end{table}
